\documentclass[12pt,a4paper,fontset=none]{ctexart}

\usepackage[margin=2.6cm]{geometry}
\usepackage{amsmath,amssymb}
\usepackage{array}
\usepackage{booktabs}
\usepackage{caption}
\usepackage{enumitem}
\usepackage{graphicx}
\usepackage{listings}
\usepackage{setspace}
\usepackage{titlesec}
\usepackage{tikz}
\usepackage{xcolor}

\usetikzlibrary{arrows.meta,positioning}
\definecolor{sysugreen}{RGB}{0,82,49}
\definecolor{codegray}{RGB}{245,245,245}
\definecolor{codekeyword}{RGB}{0,92,175}
\definecolor{codecomment}{RGB}{0,128,88}
\definecolor{codestring}{RGB}{172,92,0}
\definecolor{codebuiltin}{RGB}{128,64,160}
% 优先使用原模板字体；未安装时使用 TeX 自带的宋体。
\IfFontExistsTF{Songti SC}{%
  \setCJKmainfont{Songti SC}
  \setCJKsansfont{Songti SC}
  \setCJKmonofont{Songti SC}
}{%
  \setCJKmainfont{FandolSong-Regular}
  \setCJKsansfont{FandolSong-Regular}
  \setCJKmonofont{FandolSong-Regular}
}
\titleformat{\section}{\normalfont\Large}{\thesection}{1em}{}
\titleformat{\subsection}{\normalfont\large}{\thesubsection}{1em}{}
\titleformat{\subsubsection}{\normalfont\normalsize}{\thesubsubsection}{1em}{}

% 允许少量额外伸缩，避免中西文混排出现轻微的右边界溢出。
\setlength{\emergencystretch}{2em}

\lstset{
  backgroundcolor=\color{codegray},
  basicstyle=\ttfamily\small,
  breaklines=true,
  columns=fullflexible,
  commentstyle=\color{codecomment},
  frame=single,
  keepspaces=true,
  keywordstyle=\color{codekeyword},
  numbers=left,
  numberstyle=\tiny\color{gray},
  showstringspaces=false
}

\lstdefinestyle{cppstyle}{
  language=C++,
  keywordstyle=\color{codekeyword},
  commentstyle=\color{codecomment},
  stringstyle=\color{codestring},
  morekeywords={mpz_class,std,filesystem,string,ostream}
}

\lstdefinestyle{bashstyle}{
  language=bash,
  keywordstyle=\color{codekeyword},
  commentstyle=\color{codecomment},
  stringstyle=\color{codestring},
  morekeywords={make,xelatex,brew,sudo,apt,install}
}

\renewcommand{\lstlistingname}{代码}

\lstdefinestyle{pythonstyle}{
  language=Python,
  keywordstyle=\color{codekeyword},
  commentstyle=\color{codecomment},
  stringstyle=\color{codestring},
  emph={np,Image,Path,PCA,loadmat},
  emphstyle=\color{codebuiltin}
}

% ====== 封面信息：按需修改这里 ======
\newcommand{\headerimage}{sysu.png}
\newcommand{\reporttitle}{Assignment 1 实验报告}
\newcommand{\reportsubtitle}{MDP}
\newcommand{\studentname}{钟泓瑨}
\newcommand{\studentid}{23336338}
\newcommand{\classid}{202610271}
\newcommand{\major}{计算机科学与技术}
\newcommand{\college}{中山大学}
\newcommand{\reportdate}{\today}
% ==================================

\newcommand{\coverfield}[2]{%
  \makebox[3.1cm][c]{#1}：&\makebox[7cm][l]{#2}\\[0.55em]
}

\newcommand{\makecover}{%
  \begin{titlepage}
    \thispagestyle{plain}
    \centering

    \vspace*{1.2cm}

    \IfFileExists{\headerimage}{%
      \includegraphics[width=0.86\textwidth]{\headerimage}%
    }{%
      {\fontsize{31}{38}\selectfont\textcolor{sysugreen}{中山大学}}\par
      \vspace{0.2cm}
      {\fontsize{15}{18}\selectfont\textcolor{sysugreen}{SUN YAT-SEN UNIVERSITY}}\par
    }

    \vspace{4.1cm}

    {\fontsize{26}{34}\selectfont \reporttitle\par}

    \vspace{1.1cm}

    {\fontsize{18}{24}\selectfont \reportsubtitle\par}

    \vfill

    \begin{spacing}{1.2}
      \large
      \begin{tabular}{>{\raggedleft\arraybackslash}p{3.6cm}l}
        \coverfield{姓名}{\studentname}
        \coverfield{学号}{\studentid}
        \coverfield{教学班号}{\classid}
        \coverfield{专业}{\major}
        \coverfield{院系}{\college}
      \end{tabular}
    \end{spacing}

    \vspace{1.7cm}

    {\large \reportdate\par}

    \vspace*{0.4cm}
  \end{titlepage}
}

\begin{document}

\makecover
\tableofcontents
\newpage

% 在此填写实验报告正文。
\section{实验要求及问题定义}
\begin{itemize}
  \item 将给定迷宫建模为马尔可夫决策过程（MDP），使用 Policy Iteration 或 Value Iteration 求解；本实验采用 Value Iteration。
  \item 状态：用智能体所在格子的坐标 $(r,c)$ 表示，墙体不属于状态空间；\texttt{S} 和 \texttt{G} 分别为起点和终点。
  \item 动作：向北、东、南、西（N、E、S、W）四个方向移动，每次一格。
  \item 状态转移：移动结果确定；若目标位置为墙体或超出边界，则留在原地。
  \item 奖励与终止：每步奖励为 $-1$，包括撞墙和进入终点；到达终点后结束，终点价值为 $0$。
  \item 目标：在折扣因子 $0<\gamma<1$ 下，求使累计折扣回报最大的策略，即以最少步数从起点到达终点。
\end{itemize}

\section{MDP建模}

将迷宫问题建模为马尔可夫决策过程 $\mathcal{M}=(\mathcal{S},\mathcal{A},P,R,\gamma)$，五个分量依次为状态空间、动作空间、状态转移概率、奖励函数和折扣因子。

状态空间 $\mathcal{S}$ 由所有可通行格子的坐标 $(r,c)$ 构成，包含起点和终点，终点记为 $s_g$。行、列编号从 $0$ 开始，分别向下、向右增大。

动作以行、列位移表示，上、右、下、左四个方向对应的动作空间为
\[
\mathcal{A}=\{(-1,0),(0,1),(1,0),(0,-1)\}.
\]
用 $s+a$ 表示坐标与位移按分量相加的结果。若目标格可通行，则移至该格；撞墙或越界时留在原地。终点为吸收状态，转移函数为
\[
f(s,a)=
\begin{cases}
s_g, & s=s_g,\\
s+a, & s\ne s_g\ \text{且}\ s+a\in\mathcal{S},\\
s, & \text{其他情况}.
\end{cases}
\]
对应的转移概率为
\[
P(s'\mid s,a)=
\begin{cases}
1, & s'=f(s,a),\\
0, & \text{其他情况}.
\end{cases}
\]

非终点状态的每次动作奖励为 $-1$，包括撞墙、越界和进入终点。到达终点后回合结束，后续奖励为 $0$。奖励函数及终点价值为
\[
R(s,a)=
\begin{cases}
-1, & s\ne s_g,\\
0, & s=s_g,
\end{cases}
\qquad V(s_g)=0.
\]

策略 $\pi$ 根据当前状态选择动作，其价值函数定义为从状态 $s$ 出发的期望累计折扣回报：
\[
V^\pi(s)=\mathbb{E}_\pi\!\left[
\sum_{t=0}^{\infty}\gamma^t R(s_t,a_t)\,\middle|\,s_0=s
\right].
\]
求解目标是得到最优价值函数 $V^*(s)=\max_\pi V^\pi(s)$ 及对应策略 $\pi^*$。对于终点可达的状态，若经过 $d$ 步到达终点，回报为 $-\sum_{t=0}^{d-1}\gamma^t$。每多走一步，回报减少 $\gamma^d$，因此最大化回报等价于选择到终点的最短路径。

\section{算法设计}

\subsection{求解思路}

由于迷宫的转移与奖励模型在第 2 节中已经完全已知，求解时不需要与环境交互采样，可以直接使用动态规划。记贝尔曼最优算子为 $\mathcal{T}$，它对任意有界价值函数 $V$ 的作用定义为
\[
(\mathcal{T}V)(s)=
\begin{cases}
0, & s=s_g,\\[2pt]
\displaystyle\max_{a\in\mathcal{A}}\Bigl\{R(s,a)+\gamma V\bigl(f(s,a)\bigr)\Bigr\}, & s\ne s_g .
\end{cases}
\]
其中终点被显式地当作吸收态：备份时不再对终点自举，而是直接把价值固定为 $V(s_g)=0$，这与终点的奖励 $R(s_g,\cdot)=0$ 及零后续奖励一致。其余状态的最优性备份只涉及一个后继状态，因此期望算子退化为一次函数求值。

该算子具有两条关键性质：
\begin{itemize}[nosep]
  \item \textbf{单调性}：若价值函数逐分量满足 $V\le U$，则 $\mathcal{T}V\le\mathcal{T}U$；
  \item \textbf{压缩性}：记最大范数 $\|V\|_\infty=\max_{s\in\mathcal{S}}|V(s)|$，则
  \[
  \|\mathcal{T}V-\mathcal{T}U\|_\infty\le\gamma\|V-U\|_\infty .
  \]
\end{itemize}
压缩性的证明只用到不等式 $|\max_a x_a-\max_a y_a|\le\max_a|x_a-y_a|$。对 $s\ne s_g$ 有
\[
\bigl|(\mathcal{T}V)(s)-(\mathcal{T}U)(s)\bigr|
\le\gamma\max_{a\in\mathcal{A}}\bigl|V\bigl(f(s,a)\bigr)-U\bigl(f(s,a)\bigr)\bigr|
\le\gamma\|V-U\|_\infty ,
\]
而在 $s=s_g$ 处两者之差为 $0$，故上式对所有状态成立。

由 $\gamma<1$ 与压缩性，Banach 不动点定理保证 $\mathcal{T}$ 存在唯一不动点 $V^*=\mathcal{T}V^*$，它就是最优价值函数；并且从任意初值出发的迭代 $V_{k+1}=\mathcal{T}V_k$ 都以几何速率 $\gamma^k$ 收敛到 $V^*$。基于这一迭代格式的算法即价值迭代（Value Iteration）。

本实验采用价值迭代而不采用策略迭代，原因是两者虽然都收敛到最优策略，但价值迭代只需要一个循环，每一轮对全部状态执行一次最优性备份，每轮开销固定且实现简单；策略迭代的外层轮数通常更少，但每一轮都要把当前策略的策略评估迭代到收敛，因而引入一个内层循环和额外的停止条件。对本问题这样规模很小、转移确定的 MDP，价值迭代在实现与验证成本上更占优势。两种方法的对比见表~\ref{tab:vi-vs-pi}。

\begin{table}[htbp]
  \centering
  \caption{价值迭代与策略迭代的对比}
  \label{tab:vi-vs-pi}
  \begin{tabular}{p{2.4cm}p{5.4cm}p{5.4cm}}
    \toprule
    比较项 & 价值迭代 & 策略迭代 \\
    \midrule
    迭代结构 & 单层循环，直接迭代最优价值函数 & 外层策略改进 + 内层策略评估 \\
    每轮开销 & $O(|\mathcal{S}||\mathcal{A}|)$，固定 & 需要一次（近似）策略评估 \\
    收敛保证 & 压缩映射，误差按 $\gamma^k$ 几何衰减 & 有限步内收敛到最优策略 \\
    停止条件 & 只需要 $\max_s|\Delta V(s)|<\varepsilon$ & 需要策略稳定与评估精度两个条件 \\
    本实验 & 采用 & 作为替代方案未采用 \\
    \bottomrule
  \end{tabular}
\end{table}

\subsection{算法流程}

代码~\ref{lst:vi} 给出同步价值迭代的伪代码。输入为第 2 节建立的模型 $(\mathcal{S},\mathcal{A},f,R,\gamma)$ 与容差 $\varepsilon$，输出为最优价值函数与贪心策略。

\begin{lstlisting}[style=pythonstyle,numbers=left,caption={同步价值迭代伪代码},label={lst:vi}]
输入: 状态集 S, 动作集 A, 转移 f, 奖励 R, 折扣 gamma, 容差 eps
输出: 价值表 V, 策略 pi

V[s] <- 0                       # 初始化，终点天然为 0
repeat
    delta <- 0
    for s in S:                 # 同步备份，先写入 V_new
        if s == goal:
            V_new[s] <- 0       # 吸收态，不自举
        else:
            V_new[s] <- max{ R(s,a) + gamma * V[f(s,a)] : a in A }
        delta <- max(delta, |V_new[s] - V[s]|)
    V <- V_new
until delta < eps               # 停止准则

for s in S, s != goal:          # 策略提取
    pi[s] <- argmax{ R(s,a) + gamma * V[f(s,a)] : a in A }
return V, pi
\end{lstlisting}

流程中的几个要点如下：
\begin{enumerate}[nosep]
  \item \textbf{初始化}：所有状态（含终点）的价值初始化为 $0$。取 $V_0\equiv 0$ 使得 $V_0\ge V^*$，迭代序列从上方单调下降逼近 $V^*$，并且 $V_0(s_g)=0$ 与终点的真实价值一致。
  \item \textbf{最优性备份}：对每个非终点状态，用当前价值表计算四个动作的 $Q$ 值并取最大，得到该状态的新价值。终点不参与备份，恒为 $0$。
  \item \textbf{同步更新}：一轮之内所有状态都读取旧价值表 $V_k$，结果写入新表 $V_{k+1}$，整轮结束后再整体替换。这样每次迭代都严格对应于算子 $\mathcal{T}$ 的一次作用，记录的最大更新量 $\delta_k$ 就是真实的最大范数差 $\|V_{k+1}-V_k\|_\infty$，可以直接用于误差估计；若改为就地更新（Gauss--Seidel），实践中收敛更快，但 $\delta_k$ 不再是算子意义下的干净上界，停止准则的误差推导会变弱。
  \item \textbf{终止判断}：当一轮中所有状态的价值改变量都不超过 $\varepsilon$ 时停止迭代。
  \item \textbf{策略提取}：收敛后只提取一次策略，避免迭代中期不可靠的价值影响最终结果。
\end{enumerate}

\subsection{收敛性与停止准则}

记 $V_k$ 为第 $k$ 轮迭代后的价值表，$V_0\equiv 0$。由于所有非终点状态的单步奖励为 $-1$、终点价值为 $0$，价值始终落在区间 $\left[-\frac{1}{1-\gamma},\,0\right]$ 内，因此价值表不会溢出，$\mathcal{T}$ 的良好定义也得到了保证。

在压缩性的基础上，迭代误差有先验界
\[
\|V_k-V^*\|_\infty\le\gamma^k\|V_1-V_0\|_\infty\le\gamma^k ,
\]
其中用到 $\|V_1-V_0\|_\infty=1$：因为 $V_1(s)=-1\ (s\ne s_g)$ 而 $V_0\equiv 0$。要使误差小于 $\varepsilon$，只需迭代轮数满足 $k\ge\dfrac{\log\bigl(\varepsilon(1-\gamma)/\gamma\bigr)}{\log\gamma}$。该估计只利用了最坏情况的初始误差，通常相当保守。

实际实现使用的停止准则是后验的。设第 $k$ 轮的实测更新量为
\[
\delta_k=\|V_{k+1}-V_k\|_\infty=\max_{s\in\mathcal{S}}\bigl|V_{k+1}(s)-V_k(s)\bigr| ,
\]
则由压缩性可得标准误差界
\[
\|V_{k+1}-V^*\|_\infty\le\frac{\gamma}{1-\gamma}\,\delta_k .
\]
因此取 $\delta_k<\varepsilon$ 作为停止条件时，返回价值的误差不超过 $\dfrac{\gamma\varepsilon}{1-\gamma}$。本实验取 $\gamma=0.99$、$\varepsilon=10^{-8}$，于是价值误差上界约为 $9.9\times10^{-7}$，已经远小于迷宫问题的步数尺度，对策略的影响可以忽略。

需要指出的是，\textbf{在本题的单位步代价与零初始化下，价值迭代会在有限轮内到达精确不动点}。设 $d(s)$ 为状态 $s$ 到终点的最短步数，$D=\max_{s}d(s)$。对 $d$ 做归纳可证：任何动作的 $Q$ 值都不超过 $-1$（因为即时奖励为 $-1$ 且价值非正），而 $V^*(s)=-\sum_{t=0}^{d(s)-1}\gamma^t\le-1$，因此朝终点方向的动作为最优动作；于是经过 $k$ 轮后所有满足 $d(s)\le k$ 的状态已经取到最优值。特别地，当 $k\ge D$ 时 $V_k=V^*$，再迭代一轮得到 $\delta_{D+1}=0$，算法判定收敛并返回精确的最优价值。这就是实验中迭代轮数约为 $D+1$、且最后一次更新量恰为 $0$ 的原因（见第 4 节）。对于一般的不确定 MDP，这一有限步精确收敛的性质不再成立，价值迭代只能给出 $\varepsilon$ 意义下的近似解。

\subsection{策略提取}

收敛后，对每个非终点状态 $s$ 按下式提取贪心动作：
\[
\pi^*(s)=\arg\max_{a\in\mathcal{A}}\Bigl\{R(s,a)+\gamma V\bigl(f(s,a)\bigr)\Bigr\}.
\]
终点为吸收态，不定义动作，故策略表中只包含 $|\mathcal{S}|-1$ 个状态。

贪心策略的最优性来自策略改进定理：由 $\mathcal{T}$ 的不动点性质，对 $V^*$ 取贪心得到的策略满足 $V^{\pi^*}=V^*$。若价值只是近似收敛，记 $\|V_k-V^*\|_\infty\le\varepsilon$，则可以证明贪心策略满足
\[
\|V^{\pi_k}-V^*\|_\infty\le\frac{2\gamma\varepsilon}{1-\gamma},
\]
即策略的次优程度随价值误差线性下降；在本题的精确收敛情形下，该误差为零，返回的策略就是最优策略。

当多个动作的 $Q$ 值相同（例如在开放区域中多个方向都能以相同步数到达终点）时，$\arg\max$ 的结果取决于动作的枚举顺序。本实验采用第 2 节定义的固定顺序，即北、东、南、西
\[
(-1,0),\quad (0,1),\quad (1,0),\quad (0,-1),
\]
并规定在取值相同时选取顺序最靠前的动作，从而使策略唯一、可复现，也保证了策略图在多次运行中保持一致。

\subsection{复杂度分析}

算法每一轮对 $|\mathcal{S}|$ 个状态、每个状态 $|\mathcal{A}|=4$ 个动作各做一次常数时间的备份，因此单轮时间复杂度为 $O(|\mathcal{S}||\mathcal{A}|)$。保持整个价值表与策略表各需 $O(|\mathcal{S}|)$ 空间，转移函数通过格点坐标即时计算，不需要预先存储 $|\mathcal{S}|\times|\mathcal{A}|$ 的转移表。

总时间取决于迭代轮数 $K$：一般情形下由先验界给出 $K=O\!\left(\dfrac{\log(\varepsilon(1-\gamma)/\gamma)}{\log\gamma}\right)$，在 $\gamma=0.99$、$\varepsilon=10^{-8}$ 时约为 $2.3\times10^3$ 轮；而在本题中由于有限步精确收敛，$K\le D+1\le|\mathcal{S}|$，实测轮数为 $2\sim47$ 轮。加上一次长度不超过 $D+1$ 的策略提取与轨迹展开，总体复杂度仍为 $O(|\mathcal{S}||\mathcal{A}|)$ 量级，实验中最大规模的 15$\times$15 迷宫（102 个状态）单次求解耗时不足 5 ms。

\subsection{实现要点}

算法在 \texttt{maze.py} 与 \texttt{experiment.py} 中实现，核心迭代部分如代码~\ref{lst:core} 所示。实现上的主要考虑如下：
\begin{itemize}[nosep]
  \item \textbf{状态与动作编码}：状态用 (行, 列) 元组表示，只枚举非墙体格子，因此价值表与策略表都是 $O(|\mathcal{S}|)$ 的字典，墙体不占用任何条目；动作直接用行列位移表示，与第 2 节的 $\mathcal{A}$ 一致。
  \item \textbf{转移与奖励}：转移函数 $f$ 与奖励 $R$ 合并实现在 \texttt{Maze.step} 中，一次调用返回（后继状态, 奖励, 是否终止）三元组。若目标格是墙体或越界，后继取原状态且奖励仍为 $-1$，从而避免把撞墙当作可以无限重复的“免费动作”。
  \item \textbf{终点处理}：\texttt{step} 在终点直接返回 $(s_g,0,\text{True})$，备份时对终止情形取 $0$ 作为后续价值，与代码~\ref{lst:vi} 中 $V(s_g)=0$ 的写法完全对应。
  \item \textbf{同步更新与误差记录}：每轮先把结果写入新字典再整体替换，同时记录 $\delta_k$，既保证了与压缩映射理论的一致性，也为第 4 节的收敛曲线提供了数据。
  \item \textbf{平局与可复现性}：策略提取使用带固定顺序的 \texttt{max}（Python 的 \texttt{max} 在并列时返回第一个元素），使平局时的动作选择确定，便于复现实验与比较不同迷宫的结果。
  \item \textbf{数值安全}：$\gamma<1$ 保证价值有界，迭代不会发散；同时以 $\delta_k<\varepsilon$ 作为唯一停止条件，避免使用固定轮数而带来的精度不确定性。
\end{itemize}

\begin{lstlisting}[style=pythonstyle,numbers=left,caption={核心迭代实现（摘自 \texttt{maze.py}）},label={lst:core}]
values = dict.fromkeys(maze.states, 0.0)   # V_0 = 0
deltas = []
while True:
    updated = {                            # 同步备份
        state: max(action_value(state, move) for move in MOVES)
        for state in maze.states
    }
    deltas.append(max(abs(updated[s] - values[s])
                      for s in maze.states))
    values = updated
    if deltas[-1] < tolerance:
        break
\end{lstlisting}

\section{实验结果}

\section{总结}

\end{document}
